
\section{Variable-Depth Analysis of Ouro 1.4B}
\label{app:ouro}

We run the published Ouro 1.4B checkpoint \citep{zhu2025ouro} at inference depths $K=1$ through $K=8$ on WikiText-103 validation.
Perplexity converges by $K=3$ and remains essentially flat from $K=4$ to $K=8$ (Figure~\ref{fig:ouro_depth}).
Post-norm hidden-state norms are constant by construction, while pre-norm norms remain low because inter-loop RMSNorm resets scale after every loop.

\begin{figure}[t]
\centering
\includegraphics[width=\linewidth]{figures/ouro_variable_depth.png}
\caption{Variable-depth inference on the published Ouro 1.4B model. Left: perplexity versus inference depth. Right: hidden-state norms before and after inter-loop RMSNorm.}
\label{fig:ouro_depth}
\end{figure}

This result is compatible with our analysis.
Ouro controls recurrent scale by normalizing between loops.
Our per-loop raw and final-only-normalized configurations instead make intermediate CE losses scale-sensitive, allowing the loss to control scale directly.


% ============================================================================
